% RESUMO: substitua o parágrafo de orientação pelo resumo da sua pesquisa.
% Título, nomes, notas e palavras-chave vêm de configuracoes/metadados.tex.
%
% Regras do manual IFPI 2024 (seções 4 e 7.4) para o resumo do artigo:
%   - parágrafo ÚNICO, sem recuo na primeira linha, justificado;
%   - frases concisas e afirmativas, verbo na voz ativa e na terceira
%     pessoa do singular; a primeira frase expõe o tema e a categoria
%     do trabalho;
%   - tipo informativo: finalidade, metodologia, resultados e conclusões;
%   - de 150 a 250 palavras no artigo;
%   - sem símbolos, fórmulas, equações ou enumeração de tópicos;
%   - espaçamento simples e tamanho 12.
\begin{center}
  \textbf{\ABNTEXchapterfont\MakeUppercase{\imprimirtitulo}}
\end{center}

% Autor e orientador em linhas distintas, alinhados à margem direita, com
% nota de rodapé de identificação (manual, seção 7.4).
\begin{flushright}
  \imprimirautor\footnote{\notadoautor}

  \imprimirorientador\footnote{\notadoorientador}

  % Havendo coorientação, remova o % da linha abaixo e crie a nota
  % correspondente em configuracoes/metadados.tex.
  % \imprimircoorientador\footnote{\notadoorientador}
\end{flushright}

\begin{center}
  \textbf{RESUMO}
\end{center}

\begingroup
\SingleSpacing
\noindent Apresente aqui, em um único parágrafo, o tema da pesquisa, seu objetivo,
os procedimentos metodológicos, os principais resultados e a conclusão.
Este texto é apenas uma orientação de preenchimento e deve ser substituído
antes da entrega. O resumo do artigo tem de 150 a 250 palavras; confira com
o curso as demais exigências.
\par
\noindent\textbf{Palavras-chave:} \palavraschave
\par
\endgroup
